\chapter{أطياف الضوء ورسالة الضوء}\label{ch:g10:light-spectra}

حزمة من ضوء نجم هي العيّنة الوحيدة التي سنمسكها يومًا من ذلك النجم، وقد
تبيّن أنها عيّنة سخيّة. انشر الحزمة في طيفها --- يفعل ذلك موشور زجاجي،
ويفعله ستار من قطرات المطر --- فيصير الضوء
رسالة: خلفية قوس قزح المستمرة تعلن درجة حرارة النجم،
وشيفرة من الخطوط الدقيقة تهجّي، عنصرًا عنصرًا،
ممّ صُنع النجم. يتعلم هذا الفصل قراءة جزأي
الرسالة كليهما --- على مصابيح المختبر أولًا، ثم على الشمس.

\section{الضوء الأبيض وطيفه}

\begin{definition}[تحليل الضوء الأبيض]\label{def:g10:light-spectra:white}
حين تعبر حزمة ضيقة من ضوء الشمس موشورًا زجاجيًّا، تخرج منتشرة
في شريط مستمر من الألوان، من الأحمر إلى البنفسجي. وهذا الشريط هو
\emph{طيف}\index{طيف} الضوء. والضوء الذي يحوي كل
الألوان المرئية، كضوء الشمس أو ضوء مصباح عادي، يسمّى
\emph{ضوءًا أبيض}\index{ضوء أبيض}. أما الضوء ذو اللون النقي الواحد، الذي
لا يقدر موشور على تحليله أكثر، فهو ضوء
\emph{أحادي اللون}\index{ضوء أحادي اللون} (كحزمة الليزر
مثلًا)؛ ومزيج عدة ألوان ضوء
\emph{متعدد الألوان}\index{ضوء متعدد الألوان}. والأداة التي
تُجري التحليل وتقرأ تركيب ضوء ما هي
\emph{المِطياف}\index{مطياف}.
\end{definition}

\begin{omfigure}
\begin{tikzpicture}[scale=1.0, font=\small]
  \draw[thick, fill=blue!6] (0,0) -- (1.4,2.6) -- (2.8,0) -- cycle;
  \draw[line width=1.6pt] (-2.4,1.95) -- (0.75,1.4)
    node[pos=0.22, above] {ضوء أبيض};
  \draw[gray] (0.75,1.4) -- (2.05,1.05);
  \draw[red, thick]              (2.05,1.05) -- (5.4,0.30);
  \draw[orange, thick]           (2.05,1.05) -- (5.36,0.08);
  \draw[yellow!75!black, thick]  (2.05,1.05) -- (5.31,-0.14);
  \draw[green!65!black, thick]   (2.05,1.05) -- (5.25,-0.36);
  \draw[blue, thick]             (2.05,1.05) -- (5.17,-0.58);
  \draw[violet, thick]           (2.05,1.05) -- (5.08,-0.80);
  \node[right] at (5.4,0.30)  {أحمر (أقلّ انحرافًا)};
  \node[right] at (5.08,-0.80) {بنفسجي (أكثر انحرافًا)};
\end{tikzpicture}

{\small موشور ينشر حزمة من \omterm{def:g10:light-spectra:white}{الضوء الأبيض} في طيفها؛ وعلى
شاشة، ترسم المروحة شريط قوس قزح مستمرًّا.}
\end{omfigure}

\begin{remark}[المواشير وقطرات المطر]\label{rem:g10:light-spectra:rainbow}
يعمل الموشور لأن الزجاج يحرف الضوء البنفسجي أكثر قليلًا من الضوء
الأحمر --- أما \emph{لماذا} يفعل ذلك فيعود إلى دراسة الانكسار، في
\cref{ch:g10:refraction}؛ وهنا لا نستعمل إلا الأثر. والماء يقوم بالعمل
نفسه: فكل قطرة مطر في زخّة عابرة تعمل عمل موشور صغير، تتلقى
ضوء الشمس الأبيض وتردّه مرتَّبًا بحسب اللون. فقوس قزح \omterm{def:g10:light-spectra:white}{طيف}
الشمس، مرسومًا عبر السماء.
\end{remark}

\begin{definition}[طول الموجة، بطاقة هوية الإشعاع]\label{def:g10:light-spectra:wavelength}
يتميز كل إشعاع \omterm{def:g10:light-spectra:white}{أحادي اللون} بعدد واحد، هو
\emph{طول موجته}\index{طول الموجة} $\lambda$، ويقاس بالنانومتر
($\qty{1}{nm} = \qty{e-9}{m}$). وطول الموجة بطاقة هوية
الإشعاع: أعطِ $\lambda$ تكن قد قلت كل ما يمكن قوله
عن اللون النقي. وتستجيب العين لأطوال الموجة بين نحو
$\qty{400}{nm}$ (بنفسجي) و$\qty{800}{nm}$ (أحمر) --- وهو \emph{الطيف
المرئي}\index{الطيف المرئي}. وخلف الطرفين يستمر الإشعاع
غير مرئي لنا: \emph{فوق بنفسجي}\index{فوق بنفسجي} دون
$\qty{400}{nm}$، و\emph{تحت أحمر}\index{تحت أحمر} فوق $\qty{800}{nm}$.
أما ما يقيسه طول الموجة حقيقةً --- أي طول تموّجة واحدة من
موجة ضوئية --- فيُشرح في \cref{ch:g10:signals-and-waves}؛ وفي هذا
الفصل، العدد نفسه هو ما يهمّ.
\end{definition}

\begin{example}[قراءة أطوال الموجة]\label{ex:g10:light-spectra:classify}
مؤشر ليزر أخضر يصدر عند $\qty{532}{nm}$: داخل المجال المرئي،
في الأخضر. ومصباح زئبقي مُبيد للجراثيم يصدر عند $\qty{254}{nm}$:
\omterm{def:g10:light-spectra:wavelength}{فوق بنفسجي}، غير مرئي (وخطر على العين لأنه غير مرئي
بالذات). وثنائي مِسرى تحكم التلفاز يصدر عند $\qty{940}{nm}$:
\omterm{def:g10:light-spectra:wavelength}{تحت أحمر} --- وجّهه نحو كاميرا هاتف يرى حسّاسها قليلًا وراء
$\qty{800}{nm}$، فتظهر الومضات غير المرئية على الشاشة.
\end{example}

\section{الأطياف المستمرة ودرجة الحرارة}

\begin{definition}[الطيف المستمر]\label{def:g10:light-spectra:continuous}
المادة الكثيفة الساخنة --- شعيرة مصباح، أو فولاذ منصهر، أو سطح
نجم --- تتوهّج، ويحوي ضوءها \emph{كل} طول موجة في
مجال واسع: شريط لون متصل لا ينقصه شيء. ومثل هذا
\omterm{def:g10:light-spectra:white}{الطيف} \emph{طيف مستمر}\index{طيف مستمر}. وهو
لا يحمل أثرًا للطبيعة الكيميائية للجسم: فقضيب حديد متوهّج بالبياض
وقضيب نحاس متوهّج بالبياض يتوهّجان سواءً. أما ما يشفّره كاملًا
فهو درجة الحرارة.
\end{definition}

\begin{definition}[درجة الحرارة المطلقة]\label{def:g10:light-spectra:kelvin}
في قوانين الإشعاع، تُعدّ \emph{درجة الحرارة المطلقة}\index{درجة الحرارة
المطلقة} لا من درجة تجمّد الماء بل من
أبرد حالة ممكنة،
$\qty{-273}{\degreeCelsius}$. \omterm{def:g10:orders-of-magnitude:unit}{والوحدة} هي \emph{الكلفن}\index{كلفن}
(رمزه $\unit{K}$)، ومقدار الدرجة فيه هو نفسه:
\[
T \text{ (بوحدة } \unit{K}) = \theta \text{ (بوحدة } \unit{\degreeCelsius}) + 273 .
\]
فغرفة في $\qty{20}{\degreeCelsius}$ هي في $\qty{293}{K}$، وسطح
الشمس، قرب $\qty{5500}{\degreeCelsius}$، هو قرب
$\qty{5800}{K}$.
\end{definition}

\begin{proposition}[الأسخن أشدّ سطوعًا وأميل إلى الزرقة (قانون فين)]\label{prop:g10:light-spectra:wien}
حين ترتفع درجة حرارة جسم متوهّج كثيف:
\begin{enumerate}
  \item يشعّ أكثر عند \emph{كل} طول موجة: فيزداد \omterm{def:g10:light-spectra:white}{الطيف} كله
        سطوعًا؛
  \item وينزلق \omterm{def:g10:light-spectra:wavelength}{طول الموجة} $\lambda_{\max}$ الذي يشعّ عنده أشدّ
        إشعاع نحو جهة أطوال الموجة القصيرة (الزرقاء)،
        وفق العلاقة
        \[
        \lambda_{\max} \times T = \qty{2.90e-3}{m.K},
        \]
        حيث $T$ \omterm{def:g10:light-spectra:kelvin}{درجة الحرارة المطلقة} \omterm{def:g10:light-spectra:kelvin}{بالكلفن}.
\end{enumerate}
\end{proposition}
\admitted

\begin{remark}[من أين جاء هذا القانون]\label{rem:g10:light-spectra:wienorigin}
يلخّص قانون فين قياسات دقيقة أُجريت على الأفران والشعيرات
في أواخر القرن التاسع عشر. أما اشتقاقه فيقتضي النظرية الكمومية
للإشعاع --- بل إن هذا \omterm{def:g10:light-spectra:white}{الطيف} بالذات هو ما \emph{عجزت}
الفيزياء الكلاسيكية عن تفسيره، وهو اللغز الذي اضطرّ
بلانك إلى اختراع الكمّ. ويُجرى الاشتقاق الأمين في
مجلد السنة الثالثة.
\end{remark}

\begin{omfigure}
\begin{tikzpicture}[scale=1.0, font=\small]
  \foreach \y in {0, 1.1, 2.2}
    \foreach \xa/\xb/\c in {0/0.6/violet, 0.6/1.35/blue,
        1.35/2.4/green!75!black, 2.4/2.85/yellow, 2.85/3.45/orange,
        3.45/6/red}
      \fill[\c] (\xa,\y) rectangle (\xb,\y+0.55);
  \foreach \xa/\xb/\o in {0/0.6/0.85, 0.6/1.35/0.7, 1.35/2.4/0.5,
      2.4/2.85/0.35, 2.85/3.45/0.25, 3.45/6/0.1}
    \fill[black, opacity=\o] (\xa,0) rectangle (\xb,0.55);
  \foreach \xa/\xb/\o in {0/0.6/0.5, 0.6/1.35/0.35, 1.35/2.4/0.2,
      2.4/2.85/0.12, 2.85/3.45/0.08, 3.45/6/0.08}
    \fill[black, opacity=\o] (\xa,1.1) rectangle (\xb,1.65);
  \node[left] at (-0.15, 0.275) {$\qty{3000}{K}$};
  \node[left] at (-0.15, 1.375) {$\qty{4500}{K}$};
  \node[left] at (-0.15, 2.475) {$\qty{6000}{K}$};
  \draw[-Stealth] (0,-0.5) -- (6.7,-0.5)
    node[right] {$\lambda$ (\unit{nm})};
  \foreach \x/\l in {0/400, 1.5/500, 3/600, 4.5/700, 6/800}
    \draw (\x,-0.42) -- (\x,-0.58) node[below] {\l};
\end{tikzpicture}

{\small الشريط المرئي من \omterm{def:g10:light-spectra:continuous}{طيف مستمر} عند ثلاث درجات حرارة:
كلما برد المنبع خفت الشريط --- والطرف الأزرق يخبو
أولًا.}
\end{omfigure}

\begin{example}[أخذ درجة حرارة الشمس]\label{ex:g10:light-spectra:suntemp}
يبلغ \omterm{def:g10:light-spectra:white}{طيف} الشمس المستمر أشدّ شدّته قرب
$\lambda_{\max} = \qty{500}{nm}$. ويعطي قانون فين درجة حرارة
سطحها:
\[
T = \frac{\qty{2.90e-3}{m.K}}{\qty{5.00e-7}{m}} = \qty{5800}{K},
\]
أي نحو $\qty{5500}{\degreeCelsius}$ --- مقيسة من على بعد $150$ مليون
كيلومتر، بموشور وعملية قسمة.
\end{example}

\begin{omfigure}
\begin{tikzpicture}
\begin{axis}[omaxis, width=9.4cm, height=5.6cm,
    xlabel={$\lambda$ (\unit{nm})}, ylabel={الشدة (وحدات اعتباطية)},
    xmin=0, xmax=2150, ymin=0, ymax=9.5,
    xtick={400, 800, 1200, 1600, 2000}, ytick=\empty]
  \fill[gray, opacity=0.12] (axis cs:400,0) rectangle (axis cs:800,9.5);
  \node[gray, anchor=south] at (axis cs:600,8.4) {المرئي};
  \addplot[omThm, thick, domain=180:2050, samples=160]
    {3e16 * x^(-5) / (exp(2400/x) - 1)};
  \addplot[omDef, thick, domain=180:2050, samples=160]
    {3e16 * x^(-5) / (exp(3600/x) - 1)};
  \draw[dashed, omThm] (axis cs:480,0) -- (axis cs:480,8.0);
  \draw[dashed, omDef] (axis cs:725,0) -- (axis cs:725,1.05);
  \node[omThm, anchor=south west] at (axis cs:620,6.4)
    {نجم ساخن، $\qty{6000}{K}$};
  \node[omDef, anchor=south west] at (axis cs:950,1.3)
    {نجم بارد، $\qty{4000}{K}$};
\end{axis}
\end{tikzpicture}

{\small الشدة بدلالة \omterm{def:g10:light-spectra:wavelength}{طول الموجة} لنجمين: الأسخن
يشعّ أكثر عند كل طول موجة، وذروته $\lambda_{\max}$ (المتقطعة)
تقع عند طول موجة أقصر.}
\end{omfigure}

\section{الأطياف الخطية للغازات المثارة}

\begin{definition}[طيف الإصدار الخطي]\label{def:g10:light-spectra:emission}
الغاز تحت ضغط منخفض، إذا أثاره تفريغ كهربائي (كما في أنبوب نيون
أو مصباح شارع بالصوديوم)، يتوهّج هو أيضًا --- لكن ضوءه في \omterm{def:g10:light-spectra:white}{المِطياف}
لا يشبه قوس قزح في شيء. فلا يحضر إلا بضعة أطوال موجة
معزولة: خطوط رفيعة ساطعة على خلفية سوداء، تسمّى \emph{خطوطًا
طيفية}\index{خط طيفي}. ومثل هذا \omterm{def:g10:light-spectra:white}{الطيف} \emph{طيف إصدار
خطي}\index{طيف الإصدار}\index{طيف خطي}.
\end{definition}

\begin{example}[الهيدروجين والصوديوم]\label{ex:g10:light-spectra:hydrogen}
يصدر الهيدروجين المثار أربعة خطوط بالضبط في المرئي: خطًّا أحمر عند
$\qty{656}{nm}$، وخطًّا أزرق مخضرًّا عند $\qty{486}{nm}$، وخطين
بنفسجيين عند $\qty{434}{nm}$ و$\qty{410}{nm}$. أما الصوديوم المثار فيصدر
خطًّا واحدًا قويًّا أصفر برتقاليًّا عند $\qty{589}{nm}$ أساسًا --- ولهذا
تغرق مصابيح الشوارع القديمة أحياء بأكملها في ذلك اللون الواحد:
إذ لا يحوي ضوءها سواه تقريبًا.
\end{example}

\begin{proposition}[بصمة العناصر]\label{prop:g10:light-spectra:fingerprint}
يصدر كل عنصر كيميائي قائمته الثابتة الخاصة من الخطوط الطيفية،
وهي نفسها في كل مختبر وفي كل عصر، ولا يشترك عنصران في
القائمة نفسها. \omterm{def:g10:light-spectra:emission}{فالطيف الخطي} يعيّن إذن العنصر
المُصدِر، كما تعيّن البصمة شخصًا.
\end{proposition}
\admitted

\begin{remark}[لماذا خطوط؟]\label{rem:g10:light-spectra:whylines}
أما لماذا لا تصدر الذرة إلا أطوال موجة بعينها --- ولماذا تفضح تلك
الأطوال العنصر --- فتفسّره النظرية الكمومية للذرة، في
السنة الأخيرة من هذا المجلد: فلكل عنصر سُلّم من مستويات
الطاقة، وكل خط يعلّم قفزة واحدة بين درجتين. أما الآن فنستعمل
البصمة كما يستعملها المحققون: دون أن نسأل الأصابع كيف نمت.
\end{remark}

\section{أطياف الامتصاص}

\begin{definition}[طيف الامتصاص]\label{def:g10:light-spectra:absorption}
أرسِل \omterm{def:g10:light-spectra:white}{ضوءًا أبيض} \emph{عبر} غاز بارد تحت ضغط منخفض، ثم حلّل
ما يخرج: يعود قوس قزح المستمر إلى الظهور، لكن تخدشه خطوط قاتمة
رفيعة. فقد نزع الغاز أطوال موجة بعينها من الضوء العابر.
ومثل هذا \omterm{def:g10:light-spectra:white}{الطيف} --- خلفية مستمرة ناقصةً مجموعة من الخطوط --- هو
\emph{طيف امتصاص}\index{طيف الامتصاص}.
\end{definition}

\begin{proposition}[خطوط الإصدار وخطوط الامتصاص تتطابق]\label{prop:g10:light-spectra:kirchhoff}
يمتصّ الغاز أطوال الموجة التي يقدر على إصدارها بالضبط: فالخطوط القاتمة
في \omterm{def:g10:light-spectra:absorption}{طيف امتصاص} عنصر تقع عند أطوال الموجة نفسها التي تقع عندها
الخطوط الساطعة في \omterm{def:g10:light-spectra:white}{طيف} إصداره. ولذلك يمكن قراءة بصمة
\cref{prop:g10:light-spectra:fingerprint} في الظلام كما تُقرأ في
الضوء.
\end{proposition}
\admitted

\begin{remark}[آلية واحدة بإشارتين]\label{rem:g10:light-spectra:mechanism}
ليس التطابق مصادفة: فالامتصاص عند $\lambda$ يصعد درجة الطاقة
نفسها التي ينزلها الإصدار عند $\lambda$ --- وفصول الكمّ
في آخر هذا المجلد تضبط ذلك. وقد أرساه كيرشوف
وبنزن سنة 1859، وهو المفتاح الذي يفتح القسم التالي كله.
\end{remark}

\begin{omfigure}
\begin{tikzpicture}[scale=1.0, font=\small]
  \fill[black] (0,1.5) rectangle (6,2.1);
  \draw[violet, line width=1.6pt] (0.15,1.5) -- (0.15,2.1);
  \draw[blue, line width=1.6pt]   (0.51,1.5) -- (0.51,2.1);
  \draw[cyan, line width=1.6pt]   (1.29,1.5) -- (1.29,2.1);
  \draw[red, line width=1.6pt]    (3.84,1.5) -- (3.84,2.1);
  \node[left] at (-0.15,1.8) {إصدار};
  \foreach \xa/\xb/\c in {0/0.6/violet, 0.6/1.35/blue,
      1.35/2.4/green!75!black, 2.4/2.85/yellow, 2.85/3.45/orange,
      3.45/6/red}
    \fill[\c] (\xa,0) rectangle (\xb,0.6);
  \foreach \x in {0.15, 0.51, 1.29, 3.84}
    \draw[black, line width=1.6pt] (\x,0) -- (\x,0.6);
  \node[left] at (-0.15,0.3) {امتصاص};
  \foreach \x in {0.15, 0.51, 1.29, 3.84}
    \draw[dashed, gray] (\x,0.6) -- (\x,1.5);
  \draw[-Stealth] (0,-0.5) -- (6.7,-0.5)
    node[right] {$\lambda$ (\unit{nm})};
  \foreach \x/\l in {0/400, 1.5/500, 3/600, 4.5/700, 6/800}
    \draw (\x,-0.42) -- (\x,-0.58) node[below] {\l};
\end{tikzpicture}

{\small طيفا الهيدروجين: خطوط الإصدار الساطعة (أعلى) وخطوط
الامتصاص القاتمة (أسفل) عند أطوال الموجة نفسها بالضبط، $410$ و$434$
و$486$ و$\qty{656}{nm}$.}
\end{omfigure}

\begin{example}[دورق من بخار الصوديوم]\label{ex:g10:light-spectra:sodiumflask}
يعبر \omterm{def:g10:light-spectra:white}{ضوء أبيض} دورقًا فيه بخار صوديوم بارد. فيكون \omterm{def:g10:light-spectra:white}{الطيف}
النافذ قوس قزح كاملًا فيه خط قاتم واحد عند $\qty{589}{nm}$ ---
وهو بالضبط حيث يسطع مصباح الصوديوم أشدّ سطوع. وبالنظر إلى الدورق من
جانبه في غرفة مظلمة، يُرى الضوء المسروق وقد أُعيد إصداره: بريق
أصفر برتقالي خافت عند \omterm{def:g10:light-spectra:wavelength}{طول الموجة} نفسه.
\end{example}

\section{التحليل الطيفي: قراءة نجم}

\begin{definition}[التحليل الطيفي]\label{def:g10:light-spectra:analysis}
\emph{التحليل الطيفي}\index{التحليل الطيفي} هو تعيين
العناصر الكيميائية التي يحويها منبع (لهب أو مصباح أو نجم)
بمطابقة خطوط طيفه على قوائم العناصر
المخبرية.
\end{definition}

\begin{method}[قراءة طيف نجمي]\label{met:g10:light-spectra:reading}
النجم كرة من مادة كثيفة ساخنة (السطح) يلفّها غلاف جوي أرقّ
وأبرد. ولذلك يصل ضوءه طيفًا مستمرًّا
تخدشه خطوط امتصاص --- ويُقرأ كل جزء على حدة:
\begin{enumerate}
  \item عيّن الذروة $\lambda_{\max}$ للخلفية المستمرة؛
        فيعطي قانون فين (\cref{prop:g10:light-spectra:wien}) درجة حرارة
        السطح $T = \qty{2.90e-3}{m.K} / \lambda_{\max}$.
  \item أحصِ أطوال موجة الخطوط القاتمة.
  \item طابقها على القوائم: ولا يُعلن عن حضور عنصر
        إلا إذا ظهرت \emph{كل} خطوطه القوية.
  \item استنتج: الامتصاص يحدث في الغلاف الجوي للنجم، فالعناصر
        المعيَّنة هي إذن عناصر الغلاف الجوي.
\end{enumerate}
\end{method}

\begin{example}[وصفة الشمس]\label{ex:g10:light-spectra:sunrecipe}
يحمل \omterm{def:g10:light-spectra:white}{الطيف} الشمسي آلاف الخطوط القاتمة، وقد رسم خريطتها فراونهوفر
سنة 1814. ومن أقواها: $410$ و$434$ و$486$ و$\qty{656}{nm}$ ---
أي بصمة الهيدروجين كاملة --- إلى جانب خط الصوديوم عند
$\qty{589}{nm}$ وزوج من الكالسيوم عند $393$ و$\qty{397}{nm}$. فيحوي
الغلاف الجوي للشمس الهيدروجين قبل كل شيء، مع آثار من عناصر
مألوفة في صخور الأرض. وتنقّح القياسات الحديثة العنوان: فالشمس
ثلاثة أرباعها تقريبًا هيدروجين بالكتلة.
\end{example}

\begin{remark}[العنصر الذي وُجد في السماء أولًا]\label{rem:g10:light-spectra:helium}
خلال كسوف سنة 1868 ظهر خط ساطع عند $\qty{587.6}{nm}$
في \omterm{def:g10:light-spectra:white}{طيف} الغلاف الجوي للشمس --- ولم يوافق أي عنصر معروف. وبحسب
\cref{prop:g10:light-spectra:fingerprint} فإن البصمة المجهولة تعني عنصرًا
مجهولًا، فأُعلن عن عنصر وسُمّي باسم شمس الإغريق،
\emph{هيليوس}: الهيليوم. ولم يُعزل على الأرض إلا سنة 1895 ---
وهو العنصر الكيميائي الوحيد الذي اكتُشف في الفضاء قبل أن يوجد في البيت.
\end{remark}

\section{تمارين}

\begin{exercise}[$\star$]\label{exo:g10:light-spectra:1}
لكل إشعاع مما يلي، قل أهو \omterm{def:g10:light-spectra:wavelength}{فوق بنفسجي} أم مرئي أم \omterm{def:g10:light-spectra:wavelength}{تحت أحمر}،
وأعطِ لونه إن كان مرئيًّا: $\qty{254}{nm}$ (مصباح زئبقي)،
$\qty{480}{nm}$، $\qty{589}{nm}$، $\qty{656}{nm}$، $\qty{1064}{nm}$
(ليزر قطع).
\end{exercise}

\begin{exercise}[$\star$]\label{exo:g10:light-spectra:2}
تسقط حزمة ضيقة من ضوء الشمس على موشور زجاجي.
\begin{enumerate}
  \item صِف ما يظهر على شاشة موضوعة خلف الموشور.
  \item أيّ طرفي الشريط انحرف أكثر؟
  \item يُستبدل بضوء الشمس حزمة مؤشر ليزر أخضر.
        فماذا يظهر على الشاشة الآن، وماذا يبيّن ذلك عن
        ضوء الليزر؟
\end{enumerate}
\end{exercise}

\begin{exercise}[$\star$]\label{exo:g10:light-spectra:3}
اربط كل منبع بنوع \omterm{def:g10:light-spectra:white}{الطيف} الذي ينتجه (مستمر، أو إصدار
خطي، أو امتصاص): (أ) شعيرة مصباح متوهّج؛
(ب) أنبوب نيون إعلاني؛ (ج) حديد منصهر في مسبك؛ (د) مصباح شارع
بالصوديوم؛ (ه) ضوء الشمس المستقبَل على الأرض.
\end{exercise}

\begin{exercise}[$\star$]\label{exo:g10:light-spectra:4}
عن أقواس قزح:
\begin{enumerate}
  \item ما الذي يقوم بدور الموشور؟
  \item لا يُرى قوس قزح إلا والشمس خلف الراصد والمطر
        أمامه. فماذا يوفّر كل مكوّن من المكوّنين؟
  \item اشرح لماذا يبرهن قوس قزح على أن ضوء الشمس \omterm{def:g10:light-spectra:white}{ضوء أبيض}.
\end{enumerate}
\end{exercise}

\begin{exercise}[$\star$]\label{exo:g10:light-spectra:5}
يسخّن حدّاد مسمارًا من حديد: فيتوهّج أحمر باهتًا، ثم برتقاليًّا زاهيًا،
ثم أبيض تقريبًا.
\begin{enumerate}
  \item رتّب المراحل الثلاث بحسب درجة الحرارة.
  \item أيّ أثرين من آثار \cref{prop:g10:light-spectra:wien} يوضّحهما
        هذا التسلسل؟
  \item برّر عبارة ``المتوهّج بالبياض أسخن من المتوهّج بالحمرة''.
\end{enumerate}
\end{exercise}

\begin{exercise}[$\star\star$]\label{exo:g10:light-spectra:6}
يبلغ \omterm{def:g10:light-spectra:white}{طيف} الشمس المستمر ذروته عند $\lambda_{\max} = \qty{500}{nm}$.
\begin{enumerate}
  \item احسب درجة حرارة سطح الشمس \omterm{def:g10:light-spectra:kelvin}{بالكلفن}.
  \item حوّلها إلى الدرجات المئوية.
  \item تحقّق من أن الذروة تقع داخل المجال المرئي، وحدّد
        لونها.
\end{enumerate}
\end{exercise}

\begin{exercise}[$\star\star$]\label{exo:g10:light-spectra:7}
يتوهّج لهب شمعة عند نحو $\qty{1800}{K}$؛ وتستقرّ صفيحة تسخين كهربائية مضبوطة
على ``دافئ'' عند نحو $\qty{500}{K}$.
\begin{enumerate}
  \item احسب $\lambda_{\max}$ لكل منبع، وأعطِ مجاله.
  \item الصفيحة ساخنة بلا شك --- فاليد الممدودة فوقها تحسّ
        الإشعاع --- ومع ذلك تبدو سوداء. اشرح.
  \item وذروة الشمعة أيضًا خارج المجال المرئي. فلماذا نرى
        اللهب أصلًا؟
\end{enumerate}
\end{exercise}

\begin{exercise}[$\star\star$]\label{exo:g10:light-spectra:8}
تعطي القوائم المخبرية، في المرئي: الهيدروجين
$\{410, 434, 486, 656\}\,\unit{nm}$؛ الصوديوم $\{589\}\,\unit{nm}$؛ الهيليوم
$\{447, 502, 588, 668\}\,\unit{nm}$. ويُحلَّل مصباحا تفريغ.
يُظهر المصباح A خطوطًا عند $434$ و$486$ و$\qty{656}{nm}$، مع
خط بنفسجي خافت قرب حدّ الرؤية؛ ويُظهر المصباح B خطًّا واحدًا قويًّا
عند $\qty{589}{nm}$. عيّن الغاز في كل مصباح، مع تبرير
دقيق لكون المصباح B لا يحوي هيليومًا.
\end{exercise}

\begin{exercise}[$\star\star$]\label{exo:g10:light-spectra:9}
يُرسَل \omterm{def:g10:light-spectra:white}{ضوء أبيض} عبر دورق من بخار صوديوم بارد.
\begin{enumerate}
  \item صِف \omterm{def:g10:light-spectra:white}{طيف} الضوء النافذ.
  \item عند أي طول موجة يقع خطه القاتم، ولماذا هناك
        بالضبط؟
  \item يُطفأ المصباح، وتُظلم الغرفة، ويُثار البخار
        الآن بتفريغ. صِف \omterm{def:g10:light-spectra:white}{الطيف} الجديد.
\end{enumerate}
\end{exercise}

\begin{exercise}[$\star\star$]\label{exo:g10:light-spectra:10}
يُظهر \omterm{def:g10:light-spectra:white}{طيف} عالي التحليل للشمس خطوطًا قاتمة عند $410$ و$434$
و$486$ و$517$ و$518$ و$\qty{656}{nm}$. القوائم: الهيدروجين
$\{410, 434, 486, 656\}\,\unit{nm}$؛ المغنيزيوم $\{517, 518\}\,\unit{nm}$.
\begin{enumerate}
  \item أيّ العناصر تكشفها هذه الخطوط؟
  \item في أي جزء من الشمس تُنتج الخطوط، ولماذا هي
        قاتمة لا ساطعة؟
\end{enumerate}
\end{exercise}

\begin{exercise}[$\star\star$]\label{exo:g10:light-spectra:11}
درجة حرارة سطح جسمك نحو $\qty{37}{\degreeCelsius}$.
\begin{enumerate}
  \item حوّلها إلى \omterm{def:g10:light-spectra:kelvin}{الكلفن} واحسب $\lambda_{\max}$ لإشعاعك
        الحراري أنت.
  \item في أي مجال يقع؟
  \item استنتج لماذا لا يتوهّج أحد على نحو مرئي في غرفة مظلمة، وكيف
        تعثر الكاميرا الحرارية على الناس مع ذلك.
\end{enumerate}
\end{exercise}

\begin{exercise}[$\star\star\star$]\label{exo:g10:light-spectra:12}
تشتغل شعيرة التنغستن في مصباح متوهّج عند $\qty{2700}{K}$.
\begin{enumerate}
  \item احسب $\lambda_{\max}$ وأعطِ مجاله.
  \item اشرح لماذا تصنع هذه المصابيح مصادر إضاءة رديئة ومدافئ
        صغيرة ممتازة.
  \item أيّ درجة حرارة للشعيرة تضع $\lambda_{\max}$ عند
        $\qty{550}{nm}$، أي في وسط المرئي؟ ينصهر التنغستن عند
        $\qty{3700}{K}$؛ استنتج لماذا لا يقدر أي مصباح بشعيرة على محاكاة
        ضوء النهار، ولماذا انتقلت الإضاءة إلى تقنيات أخرى.
\end{enumerate}
\end{exercise}

\begin{exercise}[$\star\star\star$]\label{exo:g10:light-spectra:13}
يُظهر \omterm{def:g10:light-spectra:white}{طيف} نجم خلفية مستمرة ذروتها عند
$\qty{420}{nm}$، تخدشها خطوط قاتمة عند $393$ و$397$ و$410$ و$434$
و$486$ و$\qty{656}{nm}$. القوائم: الهيدروجين
$\{410, 434, 486, 656\}\,\unit{nm}$؛ الكالسيوم $\{393, 397\}\,\unit{nm}$؛
الصوديوم $\{589\}\,\unit{nm}$.
\begin{enumerate}
  \item احسب درجة حرارة سطح النجم. أهو أسخن من الشمس أم
        أبرد؟
  \item أيّ العناصر يحويها غلافه الجوي يقينًا؟
  \item أوفير الصوديوم في هذا الغلاف الجوي؟ برّر.
\end{enumerate}
\end{exercise}

\begin{exercise}[$\star\star\star$]\label{exo:g10:light-spectra:14}
يحوي مصباح تفريغ مزيجًا من غازين. ويُظهر طيفه
خطوطًا عند $410$ و$434$ و$447$ و$486$ و$502$ و$588$ و$656$
و$\qty{668}{nm}$. باستعمال قوائم \cref{exo:g10:light-spectra:8}:
\begin{enumerate}
  \item عيّن الغازين.
  \item يزعم زميل لك أن الخط $\qty{588}{nm}$ يبرهن على أن المصباح
        يحوي صوديومًا أيضًا (خطه عند $\qty{589}{nm}$)، لأن
        مِطيافًا رخيصًا لا يقدر على فصل طولي موجة بينهما $\qty{1}{nm}$.
        فمن دون أداة أفضل، لماذا يكفي الهيليوم
        تفسيرًا لهذا الخط --- وأيّ قياس يحسم
        مسألة الصوديوم نهائيًّا؟
\end{enumerate}
\end{exercise}

\begin{exercise}[$\star\star\star$]\label{exo:g10:light-spectra:15}
في كوكبة الجبار، يتوهّج إبط الجوزاء أحمر بيّنًا ويتوهّج رجل الجبار
أبيض مزرقًّا. وتبلغ أطيافهما المستمرة ذروتها عند نحو $\qty{850}{nm}$
و$\qty{240}{nm}$ على الترتيب.
\begin{enumerate}
  \item احسب درجتي حرارة السطحين.
  \item استخرج نسبة درجتي الحرارة مباشرة من طولي
        الموجة، دون إعادة حساب أي منهما.
  \item اشرح اللونين، مع أن \emph{كلتا} الذروتين تقع
        خارج المجال المرئي.
  \item النجمان كلاهما ساطع للعين المجردة. فلماذا لا تجعل الذروةُ
        الواقعة خارج المرئي نجمًا غير مرئي؟
\end{enumerate}
\end{exercise}

\section{مسألة: قراءة ضوء النجوم}

\begin{problem}[{مسألة نهاية الأسبوع --- قراءة ضوء النجوم: الفيلسوف
الذي قال ``أبدًا''، وكيف يأخذ موشور درجة حرارة نجم ويقرأ
وصفته الكيميائية}]
\label{pb:g10:light-spectra:1}
سنة 1835 احتاج الفيلسوف أوغست كونت إلى مثال على معرفة
تستعصي على البشر إلى الأبد، فاختاره بعناية: التركيب
الكيميائي للنجوم. فلن يعبّئ أحد نجمًا في قارورة يومًا. وسنة 1859
طابق كيرشوف وبنزن الخطوط القاتمة في الشمس على الخطوط
الساطعة في ألهبة المختبر، فوصلت المعرفة المستعصية
بعودة الضوء. وتعيد هذه المسألة السير في الطريق كلها: بصمات
المختبر، والأجسام المتوهّجة، وخطوط الشمس القاتمة --- وأخيرًا
نجم ستحدّد درجة حرارته وتركيبه وأنت على
مكتبك.

\medskip
\noindent\textbf{الجزء I --- بصمات في السجل.}
\begin{enumerate}
  \item أعطِ مجال أطوال الموجة المرئية بالنانومتر واللون
        عند كل طرف. وأين يقع $\qty{550}{nm}$؟
  \item يُظهر مصباح تفريغ بالهيدروجين خطوطًا عند $410$ و$434$ و$486$
        و$\qty{656}{nm}$. أعطِ لون كل خط.
  \item يصدر مصباح صوديوم خطًّا واحدًا أساسًا عند $\qty{589}{nm}$.
        فكيف يبدو المصباح للعين، ولماذا يستحيل
        الحكم على لون سيارة واقفة تحت مصباح شارع كهذا؟
  \item اشرح لماذا تعيّن مجموعة من الخطوط الطيفية عنصرًا ---
        ولماذا يجب أن تُوجد \emph{كل} الخطوط القوية للعنصر
        قبل الإعلان عن حضوره.
  \item تنبّأ \omterm{def:g10:light-spectra:wavelength}{بالطيف المرئي} (أ) حين يعبر \omterm{def:g10:light-spectra:white}{ضوء أبيض} دورقًا
        من بخار صوديوم بارد، و(ب) حين يُثار البخار نفسه، وحده في
        غرفة مظلمة، بتفريغ.
\end{enumerate}

\medskip
\noindent\textbf{الجزء II --- الأجسام الساخنة وألوانها.}
\begin{enumerate}[resume]
  \item اذكر أثري رفع درجة حرارة جسم متوهّج
        كثيف في طيفه المستمر.
  \item تشتغل شعيرة هالوجينية عند $\qty{3400}{K}$. احسب
        $\lambda_{\max}$ وأعطِ مجاله.
  \item يبلغ \omterm{def:g10:light-spectra:white}{طيف} الشمس المستمر ذروته عند $\qty{500}{nm}$. احسب
        درجة حرارة سطحها.
  \item تقع تلك الذروة في الأخضر --- ومع ذلك لا تبدو الشمس
        خضراء. اشرح.
  \item قضيب متروك في الكور يُحسّ أولًا من بعيد دون
        أن يُرى متوهّجًا، ثم يتوهّج أحمر باهتًا، ثم برتقاليًّا مبيضًّا.
        اشرح التسلسل كله بالسؤال 6.
\end{enumerate}

\medskip
\noindent\textbf{الجزء III --- خطوط الشمس القاتمة.}
\begin{enumerate}[resume]
  \item \omterm{def:g10:light-spectra:white}{الطيف} الشمسي شريط مستمر تعبره آلاف
        الخطوط القاتمة الدقيقة. فأيّ جزء من الشمس ينتج الخلفية
        المستمرة، وأيّ جزء ينتج الخطوط؟
  \item تقع خطوط شمسية قوية عند $393$ و$397$ و$410$ و$434$ و$486$
        و$589$ و$\qty{656}{nm}$. القوائم: الهيدروجين
        $\{410, 434, 486, 656\}\,\unit{nm}$؛ الصوديوم $\{589\}\,\unit{nm}$؛
        الكالسيوم $\{393, 397\}\,\unit{nm}$. فأيّ العناصر حاضرة في
        الغلاف الجوي للشمس؟
  \item خطوط الهيدروجين أقوى الخطوط بفارق كبير. فماذا يوحي
        ذلك عن تركيب الشمس --- وماذا تقول القياسات
        الحديثة؟
  \item سنة 1868 لم يوافق خط عند $\qty{587.6}{nm}$ رُصد في الغلاف الجوي
        للشمس أي قائمة. أعِد بناء الاستدلال الذي
        برّر الإعلان عن عنصر جديد، وقل كيف انتهت الحكاية.
  \item لماذا تكون الخطوط الشمسية قاتمة، بينما تعطي الذرات نفسها في
        تفريغ مخبري خطوطًا ساطعة؟ وخلال كسوف كلي،
        لبضع ثوانٍ، لا يبقى مرئيًّا أمام سماء مظلمة إلا غلاف الشمس
        الرقيق --- فتومض الخطوط نفسها
        \emph{ساطعة}. اشرح.
\end{enumerate}

\medskip
\noindent\textbf{الجزء IV --- نجم على المكتب.}
يوضع أمامك \omterm{def:g10:light-spectra:white}{طيف} النجم الساطع النسر الواقع: خلفية
مستمرة ذروتها عند $\qty{290}{nm}$، تخدشها خطوط قاتمة قوية عند
$410$ و$434$ و$486$ و$\qty{656}{nm}$.
\begin{enumerate}[resume]
  \item احسب درجة حرارة سطح النسر الواقع.
  \item تقع الذروة في \omterm{def:g10:light-spectra:wavelength}{فوق البنفسجي}. تنبّأ بلون النجم
        الظاهري، ووفّق بين ``ذروة خارج المرئي'' و``نجم
        ساطع''.
  \item عيّن العنصر المهيمن على الغلاف الجوي للنسر الواقع.
  \item قارن النسر الواقع بالشمس: نسبة درجتي الحرارة، واللون،
        والسطوع لكل \omterm{def:g10:orders-of-magnitude:unit}{وحدة} من السطح المتوهّج.
  \item الخاتمة --- املأ بطاقة الهوية التي أعلن كونت
        استحالتها: درجة الحرارة، والعنصر المهيمن، واللون الظاهري
        للنسر الواقع، وكلها مقروءة من حزمة ضوء واحدة. فكم سنة
        دامت كلمة ``أبدًا''؟
\end{enumerate}
\end{problem}
